%%%%%%%%%%%%%%%%%%%%%%%%%%%%%%%%%%%%%%%%%
%%% WP01
%%%%%%%%%%%%%%%%%%%%%%%%%%%%%%%%%%%%%%%%%

\begin{prtwo}
\tsubsubsection{WP1 - Project Management and Coordination}
\end{prtwo}


%%%%%%%%%%%%%%%%%%%%%%%%%%%%%%%%%%%%%%%%%
%%% Section content, please change!
%%%%%%%%%%%%%%%%%%%%%%%%%%%%%%%%%%%%%%%%%

\subsubsection*{Overview and Goals}

\begin{table}[H]
    \renewcommand{\arraystretch}{1.50}		
    \footnotesize   
    {\color{black}\begin{tabular}{*{3}{|p{0.10\textwidth}}|l|}
        \hline
        \rowcolor{mygray} \multicolumn{4}{|c|}{\textit{\color{white}Work Package Summary}} \\
        \hline
        \rowcolor{mylightergray} \textit{WP No.} & \cellcolor{white} 1 & \textit{Title of WP} & \cellcolor{white} Project Management and Coordination (MGT) \\
        \hline
        \rowcolor{mylightergray} \textit{Start} & \cellcolor{white} M1 & \textit{End} & \cellcolor{white} M48 \\
        \hline
        \rowcolor{mylightergray} \multicolumn{4}{|p{0.978\textwidth}|}{\textit{Participating Organisations}} \\
        \hline
        \multicolumn{4}{|p{0.978\textwidth}|}{
            \hspace*{-0.75cm} 
            \begin{minipage}[t]{\textwidth}
    			\begin{itemize}
    			    \item WP Leader: 1-INFN
    				\item Participants: INFN, GANIL, CERN
    			\end{itemize} 
    			\vspace*{0.10em}
			\end{minipage}
        } \\
        \hline
    \end{tabular}}
    \vspace{0.5em}\vfill
    {\color{black}\begin{tabular}{|p{0.978\textwidth}|}
        \hline
        \rowcolor{mylightergray} \textit{Goals} \\
        \hline
        \rowcolor{white} 
        \hspace*{-0.75cm} 
        \begin{minipage}[t]{\textwidth}
        \setlength{\parindent}{15pt}
                Task 1.1. Project management and coordination
    		\begin{itemize}
    		    \item Management and steering of the whole project
                    \item Monitoring of the scientific and technical progress in all Work Packages
                    \item Ensuring the contractual and administrative implementation
                    \item Following and reporting on the use of resources
                    \item Preparation of the periodic and final project reports
    		\end{itemize} 
    		\vspace*{0.10em}
	\end{minipage}        
        \\
        \hline
    \end{tabular}}
    \vspace{0.5em}\vfill
    \resizebox{\textwidth}{!}{%
        {\color{black}\begin{tabular}{|p{0.25\textwidth}|*{3}{>{\centering\arraybackslash}p{0.252\textwidth}|}}
                    %    {|l|*{3}{>{\centering\arraybackslash}p{0.2\textwidth}|}}
            \hline    
            \rowcolor{mylightergray} \textit{Participant number} & \textit{1} & \textit{2} & \textit{3} \\
            \hline
            \rowcolor{white} \cellcolor{mylightergray}\textit{Participant short name} & INFN & GANIL & CERN \\
            \hline
            \rowcolor{white} \cellcolor{mylightergray}\textit{PM per participant~\footnotemark} & 48 & - & 24 \\
            \hline        
        \end{tabular}}
        }
\end{table}
\footnotetext{the PM figures correspond to the full duration of the project.}




%\begin{prtwo}
\subsubsection*{Status}
The Project Office (PO), coordinated by INFN, was responsible for managing all administrative processes
of EURO-LABS. Together with the Management Team (MT) and the Steering Committee
(SC), the PO ensures the smooth and timely execution of the project and reports to the Governing
Board (GB). The PO also serves as the liaison with the EU Project Officer. There was no change in the composition of the MT.
%\end{prtwo}

\subsubsection*{Progress per Task}

\subparagraph{Task 1:} Project management and coordination \mbox{}

%\begin{prtwo}
WP1 handled administrative processes and communications with REA and all the beneficiaries. The PO along with the Steering Committee, ensured a smooth running of the project. Two in-person an-nual meetings further strengthened the ties between the communities. A large focus on the presentation of selected scientific and technological results was adopted. Talks on the European Strategy of the three communities highlighted the role of EURO-LABS. All the Deliverables and Milestones have been completed in P3 and are available through open access. The reports are available through the website (https://web.infn.it/EURO-LABS/results/). The participants were given support for the implementation of the project, specifically in TNAs and financial aspects, through constant contacts and daily helpdesk via email, telephone and remote meetings.
%\end{prtwo}

%\begin{prtwo}
\subsubsection*{Main Results and Achievements}
The meetings of the Steering Committee/ Management Team were held on a regular basis (at least bi-monthly). Also periodic management meetings of the Project Office staff were held, involving people from INFN Bologna, INFN Frascati and CERN. Two annual meetings were held in the RP3 period. FAME (Fourth Annual Meeting) was held in Ljubljana, Slovenia (29th September - 1st October 2025, 57 in-person participants). The Final Annual Meeting (FINE) was organized by and hosted at KIT in Karlsruhe, Germany (29th June - 1st July 2026, 63 in-person participants). Participants also connected through remote access to the conferences. The two GB meetings held during the annual meetings were actively attended. The reports are available on the Funding and Tenders Portal and, when public, on the Zenodo repository and linked to the project website.
%\end{prtwo}


\subsubsection*{Deviations and Corrective Actions}
\label{sec:wp1_deviations}

\noindent\textbf{Grant Agreement Amendment and Budget Redistribution}

During the third and final reporting period (RP3), a major management activity involved the preparation and finalisation of a comprehensive Grant Agreement Amendment. The primary objective of this amendment was to implement an efficient budget redistribution across the consortium, aligning financial resources with the actual spending trends of the beneficiaries and responding to the demand (either higher or lower than initially foreseen) for Transnational Access (TA) to the various research infrastructures. This internal reallocation did not impact the overall grant parameters, maintaining the total consortium eligible costs at EUR 14,703,604.06 and the maximum grant amount at EUR 14,174,110.00.

Concurrently, the amendment addressed the legal transformation of Beneficiary 17 (ATOMKI) following a national institutional reform in Hungary, which required the assignment of a new PIC number. Within the Horizon Europe framework, a PIC change entails the formal addition of a new beneficiary rather than a simple administrative update. In close coordination with the European Commission officers, the most streamlined solution was adopted: Beneficiary 17 was retained in the consortium list but ceased all project activities on 28/10/2025, while the newly established entity, BEN-34 HUN-REN ATOMKI, entered the project as of 29/10/2025. Following an overall budget reduction of EUR 11,000.00 under the general redistribution, the unspent portion of BEN-17's budget (EUR 18,982.41) was transferred to BEN-34, with no costs claimed by BEN-17 after its termination date.

\noindent\textbf{WP1 Person-Month (PM) Effort Distribution}

As summarised in the table below, the overall resource consumption within WP1 (Management) at the end of the project is fully in line with the initial plan, totalling 76.19 PMs against a planned budget of 72.00 PMs (105.82\% overall execution):

{\fontsize{9}{11}\selectfont
\begin{center}
  {\color{black}\begin{tabular}[t]{!{\color{mygray}\vrule}p{0.10\linewidth}!
  {\color{mygray}\vrule}>{\centering\arraybackslash}p{0.09\linewidth}!
  {\color{mygray}\vrule}>{\centering\arraybackslash}p{0.09\linewidth}!
  {\color{mygray}\vrule}>{\centering\arraybackslash}p{0.09\linewidth}!
  {\color{mygray}\vrule}>{\centering\arraybackslash}p{0.09\linewidth}!
  {\color{mygray}\vrule}>{\centering\arraybackslash}p{0.09\linewidth}!
  {\color{mygray}\vrule}>{\centering\arraybackslash}p{0.12\linewidth}!{\color{mygray}\vrule} } \hline
    \rowcolor{mycyan} & {\bf RP1} & {\bf RP2} & {\bf RP3} & {\bf Total} & {\bf Budget} & {\bf Percentage} \\ \hline
    \cellcolor{mycyan}{\bf INFN} & 12.70 & 22.19 & 30.11 & 65.00 & 48 & 135.42\% \\ \hline
    \cellcolor{mycyan}{\bf CERN} & 2.49 & 2.17 & 6.53 & 11.19 & 24 & 46.63\% \\ \hline
    \cellcolor{mycyan}{\bf Total} & 15.19 & 24.36 & 36.64 & 76.19 & 72 & 105.82\% \\ \hline
  \end{tabular}}
\end{center}
}

As shown by the progression across the three reporting periods, INFN reached 135.42\% of its planned WP1 effort, while CERN reached 46.63\%. The higher PM consumption recorded by INFN---particularly during RP3 (30.11 PMs)---directly reflects the substantial administrative workload absorbed by the Coordinator to manage the prolonged amendment process, the complex budget redistribution across multiple partners, and the overall final reporting and closure activities, while fully compensating for the lower effort claimed by CERN within this specific work package.

\noindent\textbf{Major Financial Deviations at Project Level}

At the conclusion of the project, three main areas of financial overspending were recorded across the consortium, largely driven by the exceptional scientific uptake of the project's opportunities:

\begin{itemize}
\item \textbf{Transnational Access Unit Costs (Cost Category D.3):}\newline
  Many beneficiaries experienced an extraordinarily high demand for access to their research infrastructures, receiving a volume of high-quality project proposals aligned with the scientific scope of EURO-LABS that far exceeded initial expectations. For several beneficiaries operating under the unit cost model for TA, this virtuous scientific performance resulted in the delivery of significantly more access units than originally planned, generating an overall consortium overspending in cost category D.3 of approximately EUR 1.25 million.
\item \textbf{Personnel Costs -- CERN (approx.\ EUR 234k overspending):}\newline
  At project level, CERN's personnel costs exceeded the revised budget by 24.95\%. This deviation stems primarily from the initial phase of the project, where recruitment delays required a higher proportion of senior staff to ensure the timely start-up of TA activities. Over the course of the project, the progressive involvement of junior staff under senior supervision brought the average PM rate steadily down (from EUR 14,990 in RP1 to EUR 10,228 in RP2 and EUR 7,293 in RP3), moving closer to the planned rate. Total funded PM consumption remained very close to the revised plan, supported by 27.06 in-kind, non-funded PMs contributed by CERN across the three periods (1.96 in RP1, 8.10 in RP2, and 17.00 in RP3). Furthermore, the deviation was particularly pronounced in WP3, where the HiRadMat facility experienced exceptional demand over the last two years---delivering nearly double the access units originally planned---which required substantial additional operational and administrative effort for user coordination, registration, and reimbursement management.
\item \textbf{Personnel Costs -- CNRS (approx.\ EUR 400k overspending):}\newline
  CNRS also recorded a significant positive deviation of personnel costs amounting to approximately EUR 400k. 
  \textbf{IJCLab}:
  The personnel cost deviation reflects the strong involvement of IJCLab in the project, which required a greater contribution from its teams than initially anticipated.
The laboratory recorded in Agate-Tempo the working time of the staff involved, including permanent staff, in order to document their contribution to the project activities and support the justification of related missions.
In particular, the personnel effort reported under WP2 was mainly related to experimental work and access to the ALTO facility.
 \textbf{GANIL}:
The initial budget included only non-permanent staff costs and did not account for the involvement of permanent staff.
However, permanent staff were actively involved in the implementation of the project activities under WP1, WP2, WP3 and WP5. Their corresponding personnel costs were therefore included in the costs declared by GANIL.
RP1 adjustment :
The adjustment for RP1 concerns corrections to personnel costs declarations, as well as travel expenses that were mistakenly declared.
RP2 adjustment :
The adjustment for RP2 concerns personnel costs that were not declared during RP2. It also includes travel costs for travels made during RP2 but paid in RP3.
  

\end{itemize}

% Due to the administrative and procedural complexity of managing this legal transition alongside the consortium-wide financial reallocation, the formal finalisation of the amendment required over a year, concluding shortly after the official end date of the project. Nevertheless, because the proposed budget modifications had already been informally reviewed and endorsed by the EC Financial Officer, beneficiaries were able to operate smoothly during the final months of the project with the practical assurance that the agreed budget adjustments would be formally approved.


\subsubsection*{Milestones and Deliverables}

\begin{prtwo}
{\fontsize{9}{11}\selectfont
\begin{center}
  {\color{black}\begin{tabular}[t]{!{\color{mygray}\vrule}p{0.10\linewidth}!
  {\color{mygray}\vrule}p{0.60\linewidth}!
  {\color{mygray}\vrule}p{0.20\linewidth}!{\color{mygray}\vrule} } \hline
    \rowcolor{mycyan} & {\bf Title} & {\bf Status} \\ \hline
    \cellcolor{mycyan}{\bf D1.1}: & Periodic Report-2 (Sept 2023-Feb 2025) &  Achieved\\ \hline
   \cellcolor{mycyan}{\bf D1.5}: & Policy Brief – 3rd release  &  Achieved\\ \hline 
  \end{tabular}}
\end{center}
}
\end{prtwo}

\subsubsection*{Project Meetings}

\begin{longtable}{|L{0.18\textwidth}|C{0.48\textwidth}|C{0.12\textwidth}|C{0.10\textwidth}|}
    \caption{WP1 meetings and coordination activities during P3.}
    \label{tab:wp1-meetings}
    \\ \hline
    \rowcolor{mycyan}
    {\bf Dates} & {\bf Meeting title and link} &
     {\bf Venue} & {\bf WP} 
    \\ \hline
    \endfirsthead
    \hline
    \rowcolor{mycyan}
        {\bf Dates} & {\bf Meeting title and link} &
     {\bf Venue} & {\bf WP} 
  \\ \hline
    \endhead
    \hline
    \endfoot

03/03/2025 &
EURO-LABS Project Office meeting (INFN, CERN) &
Zoom & WP1 \\ \hline

12/03/2025 &
Meeting with GANIL-SPIRAL2 staff -- Data Reporting Support &
Zoom & WP1 \\ \hline

13/03/2025 &
\href{https://agenda.infn.it/event/45159/}
{Management Team \& Steering Committee meeting} &
Zoom & WP1 \\ \hline

19/03/2025 &
Meeting with LNL facility coordinator -- Data reporting support &
Zoom & WP1 \\ \hline

14/04/2025 &
\href{https://agenda.infn.it/event/46743/}
{Management Team \& Steering Committee meeting} &
Zoom & WP1 \\ \hline

16/04/2025 &
Meeting with WP4 Leader -- facilities TA data &
Zoom & WP1 \\ \hline

14/05/2025 &
Management Team \& Steering Committee meeting &
Zoom & WP1 \\ \hline

15/05/2025 &
\href{https://agenda.infn.it/event/47184/}
{Review Meeting (Management Team \& Steering Committee,
Project Officer REA, external reviewer)} &
Zoom & WP1 \\ \hline

22/05/2025 &
Management Team \& Steering Committee meeting &
Zoom & WP1 \\ \hline

23/05/2025 &
Management Team \& Steering Committee meeting &
Zoom & WP1 \\ \hline

26/05/2025 &
EURO-LABS Project Office meeting INFN and INFN facilities &
Teams & WP1 \\ \hline

11/06/2025 &
\href{https://cnrs.zoom.us/j/97354004837?pwd=6SEMM7rebRTsPlDO6PDfYmylhzpxQS.1}
{Management Team \& Steering Committee meeting} &
Zoom & WP1 \\ \hline

10/07/2025 &
\href{https://agenda.infn.it/event/53029/}
{Management Team \& Steering Committee meeting} &
Zoom & WP1 \\ \hline

29/09--01/10/2025 &
\href{https://indico.cern.ch/event/1572254/overview}
{Annual meeting (FAME) and GB meeting} &
JSI \& Zoom & WP1 \\ \hline

23/01/2026 &
\href{https://agenda.infn.it/event/50218/}
{Management Team \& Steering Committee meeting} &
Zoom & WP1 \\ \hline

29/01/2026 &
\href{https://cnrs.zoom.us/j/99645413746?pwd=ttYU9DUoT0MCme5cMv2aWTKnKoiWEO.1}
{Meeting with Scientific Coordinators (amendment, FINE organization)} &
Zoom & WP1 \\ \hline

05/02/2026 &
Meeting with Warrant Hub staff -- Social media and Newsletter &
Teams & WP1 \\ \hline

11/02/2026 &
\href{https://cnrs.zoom.us/j/94198661686?pwd=vXSIJGZSKqdj4bF2ezFi834t6tEmcN.1}
{Meeting with KIT staff -- organization of Final meeting FINE} &
Zoom & WP1 \\ \hline

23/02/2026 &
Meeting with Warrant Hub staff -- Social media and Newsletter &
INFN & WP1 \\ \hline

27/02/2026 &
\href{https://agenda.infn.it/event/51115/}
{Management Team \& Steering Committee meeting} &
Zoom & WP1 \\ \hline

19/03/2026 &
\href{https://agenda.infn.it/event/51115/}
{Management Team \& Steering Committee meeting} &
Zoom & WP1 \\ \hline

27/03/2026 &
Internal staff meeting (amendment) &
Teams & WP1 \\ \hline

27/04/2026 &
\href{https://agenda.infn.it/event/52026/}
{Management Team \& Steering Committee meeting} &
Zoom & WP1 \\ \hline

18/05/2026 &
Meeting with Warrant Hub staff and Deputy Scientific
Coordinator -- Dissemination Register (final report) &
Teams & WP1 \\ \hline

21/05/2026 &
\href{https://agenda.infn.it/event/52457/}
{Management Team \& Steering Committee meeting} &
Zoom & WP1 \\ \hline

29/05/2026 &
Meeting with PO &
Zoom & WP1 \\ \hline

29/06--01/07/2026 &
\href{https://indico.kit.edu/event/5341/}
{Final meeting (FINE) and GB meeting} &
KIT \& Zoom & WP1 \\ \hline

04/09/2026 &
\href{https://agenda.infn.it/event/51351/}
{Management Team \& Steering Committee meeting} &
Zoom & WP1 \\ \hline

\end{longtable}


%  {\color{mygray}\vrule}p{0.40\linewidth}!

%%%%%%%%%%%%%%%%%%%%%%%%%%%%%%%%%%%%%%%%%
